\fancyhead{}
\fancyfoot{}
\cfoot{\thepage}


\lhead{Capitulo 3}

\chapter{Herramientas Tecnológicas}

En el presente capítulo se describen las herramientas tecnológicas utilizadas para el desarrollo del proyecto.

\section{Selección de Herramientas}

\subsection{Figma}

Figma es una herramienta de diseño gráfico basada en la nube, utilizada principalmente para el diseño de interfaces de usuario (UI), experiencia de usuario (UX) y el prototipado de aplicaciones web y móviles. 

Sus características más destacadas incluyen:

\begin{itemize}
    \item Colaboración en tiempo real, permitiendo que varios usuarios trabajen simultáneamente en un mismo archivo.
    \item Funciona directamente desde el navegador, siendo multiplataforma y sin necesidad de instalación.
    \item Uso de componentes reutilizables y estilos compartidos para mantener coherencia en los diseños.
    \item Facilita la comunicación entre diseñadores y desarrolladores, permitiendo inspeccionar propiedades como CSS, medidas y exportar recursos fácilmente.
\end{itemize}

Es una herramienta ampliamente adoptada por equipos de desarrollo ágil y entornos que implementan sistemas de diseño\cite{figma}.


\subsection{Python}
Python es un lenguaje de programación de alto nivel, interpretado y de código abierto, creado por Guido van Rossum en 1991. Fue diseñado centrándose en la simplicidad y la legibilidad del código, hecho que permite a los desarrolladores expresar conceptos en menos líneas de código en comparación con otros lenguajes como C++ o Java. Descifrando el lenguaje de programación, resulta mucho más sencillo tanto de aprender como de utilizar\cite{python}.

\subsection{JavaScript}

JavaScript o JS es un lenguaje de programación diseñado para escribir las partes \textit{Front End} y \textit{Back End} de sitios web y aplicaciones móviles. Es un lenguaje de alto nivel, legible para los humanos, pero no requiere un compilador para traducirlo al lenguaje de máquina. El código se interpreta y se ejecuta directamente en un navegador web, lo que facilita y acelera la interacción con el usuario.

JavaScript es un lenguaje de \textit{scripts}, es decir, de secuencias de comandos. Permite agregar muchas funcionalidades a las páginas web haciéndolas más interactivas, útiles y atractivas para el usuario. Algunos ejemplos de contenido dinámico son menús desplegables, formularios, galerías de fotos, gráficos animados y más\cite{javascript}.




\subsection{SQL}

El lenguaje de consulta estructurado (SQL) es un lenguaje de programación estandarizado y específico de dominio que sobresale en el manejo de relaciones de datos. Se emplea ampliamente para almacenar, manipular y recuperar datos en sistemas como MySQL, SQL Server y Oracle.

Cuando es necesario recuperar datos de una base de datos, se emplea SQL para realizar la solicitud. 

SQL es utilizado por los administradores de bases de datos, desarrolladores y analistas de datos para tareas como definición de datos, control de acceso, uso compartido de datos, escritura de scripts de integración de datos y ejecución de consultas analíticas\cite{sql}.


\subsection{PostgreSQL}

Es un potente sistema de base de datos relacional de objetos de código abierto que utiliza y amplía el lenguaje SQL, combinado con muchas características que permiten almacenar y escalar de forma segura las cargas de trabajo de datos más complejas. Los orígenes de PostgreSQL se remontan a 1986, como parte del proyecto ``Postgres'' en la Universidad de California en Berkeley, y cuenta con más de 35 años de desarrollo activo en su plataforma central.

Se ha ganado una sólida reputación por su arquitectura probada, confiabilidad, integridad de datos, conjunto de características robustas, extensibilidad y por la dedicación de la comunidad de código abierto detrás del software, que ofrece soluciones innovadoras y de alto rendimiento de manera constante. PostgreSQL se ejecuta en todos los principales sistemas operativos \cite{postgresql}.

\subsection{Git}

Es un sistema de control de versiones distribuido, lo que significa que un clon local del proyecto es un repositorio de control de versiones completo. Estos repositorios locales plenamente funcionales permiten trabajar sin conexión o de forma remota con facilidad. Los desarrolladores confirman su trabajo localmente y, a continuación, sincronizan la copia del repositorio con la del servidor. Este paradigma es distinto del control de versiones centralizado, donde los clientes deben sincronizar el código con un servidor antes de crear nuevas versiones\cite{git}.



\section{Entorno de dessarrollo del modelo}

\subsection{Visual Studio Code}
Para la etapa de codificación del modelo y los componentes del sistema, se utilizó \textit{Visual Studio Code} (VS Code), un editor de código fuente desarrollado por Microsoft, diseñado para ser ligero, rápido y altamente personalizable. A diferencia de otros editores de código, como los IDE completos (Entornos de Desarrollo Integrados), VS Code se enfoca en ofrecer una experiencia ágil para la escritura y edición de código, sin sacrificar características avanzadas \cite{vscode}. Lo que lo diferencia es su capacidad de ser extendido con una gran cantidad de extensiones y herramientas, adaptándose a las necesidades específicas de cada desarrollador.

\subsection{ GitHub Desktop}
Para gestionar el control de versiones del proyecto y mantener un historial estructurado de los cambios realizados durante el desarrollo, se empleó \textit{GitHub Desktop}. Esta herramienta gratuita y de código abierto proporciona una interfaz gráfica para trabajar con repositorios Git alojados en GitHub u otros servicios. Facilita la realización de operaciones comunes de Git, como clonar, confirmar (\textit{commit}), crear ramas y sincronizar cambios, sin necesidad de utilizar la línea de comandos. Está diseñada para simplificar el flujo de trabajo tanto para principiantes como para usuarios avanzados \cite{githubdesktop}.


\subsection{Postman}
Finalmente, para la validación de los endpoints del API y la verificación de la comunicación entre los módulos del sistema, se utilizó \textit{Postman}. Es una plataforma orientada a la prueba y gestión de APIs, que facilita a los desarrolladores la verificación y documentación de servicios web. Inicialmente creada como una extensión del navegador, hoy está disponible como aplicación independiente para Windows, Linux y macOS. Su interfaz gráfica permite enviar solicitudes HTTP, automatizar pruebas y generar documentación, optimizando el desarrollo y reduciendo errores en la comunicación con servicios externos \cite{imaginaPostman}.




\section{Librerías utilizadas}

\subsection{React}

Para el desarrollo del frontend y del panel de control (\textit{dashboard}) del sistema, se utilizó \textit{React} en su versión 19.1.1. React es una biblioteca JavaScript de código abierto desarrollada por el equipo de Meta, en colaboración con una amplia comunidad de desarrolladores. Desde su lanzamiento en 2013, se ha consolidado como una de las tecnologías más utilizadas para la construcción de interfaces de usuario dinámicas y escalables.

Su funcionamiento se basa en la creación de componentes reutilizables que integran estructura y lógica mediante la sintaxis JSX. Este enfoque permite gestionar la interfaz de forma eficiente mediante un \textit{Virtual DOM}, el cual detecta cambios de estado y actualiza únicamente los elementos necesarios sin recargar toda la página, optimizando el rendimiento del sistema \cite{react}.

\subsection{Flask}
Para el desarrollo del backend y la implementación de las rutas del API del sistema, se utilizó el microframework \textit{Flask} en su versión 3.1.2. Flask es un microframework web de código abierto escrito en Python, orientado al desarrollo rápido y flexible de aplicaciones. Su estructura modular y extensible otorga a los desarrolladores un control total sobre la arquitectura, destacándose por su sencillez, ligereza e integración con múltiples extensiones.

Estas características lo convierten en una herramienta ideal para la creación de APIs RESTful y servicios que requieren agilidad y adaptabilidad, motivo por el cual fue seleccionado para la construcción de los endpoints del sistema \cite{flask}.

\subsection{YOLO}
Para la etapa de detección facial del sistema se utilizó el modelo \textit{YOLOv8}, implementado mediante la librería \textit{Ultralytics} en su versión 8.3.191. YOLO (You Only Look Once) es un algoritmo de detección de objetos en tiempo real basado en redes neuronales convolucionales. Introducido por Joseph Redmon y colaboradores en 2015, YOLO se caracteriza por procesar una imagen completa en una sola pasada, dividiéndola en una cuadrícula y prediciendo simultáneamente las clases y ubicaciones de los objetos presentes. 

Esta arquitectura unificada permite una detección rápida y precisa, siendo especialmente útil en aplicaciones que requieren procesamiento en tiempo real, motivo por el cual fue seleccionado para la detección facial dentro del sistema \cite{yolo}.


\subsection{DeepFace}
Para la etapa de reconocimiento facial se utilizó la librería \textit{DeepFace} en su versión 0.0.95. DeepFace es un framework de código abierto basado en \textit{Python} que unifica diversas arquitecturas de reconocimiento facial y facilita su implementación mediante funciones de alto nivel. Entre sus características principales se destacan el soporte para múltiples modelos de extracción de embeddings, detección facial, análisis demográfico y verificación de identidad \cite{DeepFace}.

En este proyecto, \textit{DeepFace} fue utilizada específicamente para el proceso de identificación facial, empleando el modelo \textit{ArcFace} como generador de embeddings, permitiendo un reconocimiento más robusto y confiable dentro del sistema desarrollado.

